% ============================================================
%  Method: Cross-modal Score Calibration (CRD)
% ============================================================
\subsection{CRD Overview}
\label{sec:method}
CRD is a post-hoc calibration layer. It modifies neither encoder, adds two small
heads, and returns a corrected score for a single pair. Figure~\ref{fig:overview}
shows the path: the pair is encoded, the additive heads are regressed onto the
standardised similarity, and inference combines the raw score with one subtraction.
No encoder parameter is updated, and inference never observes a label.

\subsection{Frozen Features and Score Standardisation}
For a pair $(I,T)$ we compute
\begin{equation}
  u = f_I(I), \qquad v = f_T(T), \qquad s(I,T) = \operatorname{sim}(u,v),
\end{equation}
with $\operatorname{sim}$ the cosine similarity. Because similarity scales differ
between encoders, the score is standardised with statistics estimated on the
training split alone,
\begin{equation}
  \mu = \operatorname*{mean}_{\mathcal{D}_{\text{train}}} s, \qquad
  \sigma = \operatorname*{std}_{\mathcal{D}_{\text{train}}} s + 10^{-8}, \qquad
  z = \frac{s-\mu}{\sigma}.
  \label{eq:z}
\end{equation}
Validation and test pairs use the same $\mu,\sigma$; no statistic is computed on
either split. Standardisation makes one $\lambda$ grid meaningful across backbones
with different ranges and places the target on the scale of the residual.

\subsection{Additive Score Estimation}
The estimated component is a sum of two head outputs and a bias,
\begin{equation}
  g_\theta(I,T) \;=\; a_\phi(u) + b_\psi(v) + b_0,
  \label{eq:g}
\end{equation}
where each head is a two-layer perceptron with one hidden layer of width $64$ and a
GELU nonlinearity. The additive form is retained in preference to a joint residual
network over concatenated features: a joint network can absorb part of the match term
$m(I,T)$, so the constraint is a design choice rather than a mathematical necessity. Here $g_\theta$ is neither a second similarity nor an estimate of image or text
noise; it is the component of $z$ predictable from the two sides jointly, with no
cross-side interaction term.

\subsection{Row-level Objective and Inference}
Each training pair contributes one target, its own standardised score $z$:
\begin{equation}
  \mathcal{L}(\theta) = \frac{1}{|\mathcal{D}_{\text{train}}|}
  \sum_{(I,T)} \operatorname{Huber}\!\left(g_\theta(I,T),\, z(I,T)\right).
  \label{eq:loss}
\end{equation}
Supervising on the pair's own score, rather than on its cycle mean, is substantive:
the matched and mismatched members of a cycle share the same cycle mean, so a
cycle-mean target asks the heads to predict a quantity that is constant within the
group being ranked. The encoders are frozen and only $\phi,\psi,b_0$ are optimised;
the regression target involves no clean/noisy label. Labels enter at exactly one
point, the selection of $\lambda$; the method is therefore not fully unsupervised and
is not described as such.

The calibrated score is
\begin{equation}
  z' \;=\; z \;-\; \lambda\, g_\theta(I,T),
  \label{eq:calib}
\end{equation}
with $\lambda\ge0$ selected on a validation split (for Attribute, $2{,}084$ cycles)
by maximising the AUROC of $z'$, over $\Lambda=\{0,0.1,\dots,2.0\}$. The grid
includes zero, so validation can decide that no calibration is warranted; in the
placebo runs it does. Selected values are not restricted to $[0,1]$ (they range from
$0.8$ to $1.4$ across our cells), so the operation is a scaled subtraction rather
than an exact removal. Inference is pair-local: the features are encoded once, the two heads
are evaluated, and \eqref{eq:calib} is returned, with no reference bank, partner
caption or test-time threshold. The heads hold $65{,}795$ parameters and cost
$0.2$\,ms; end to end the calibrated score adds $0.83$\,ms ($5.5\%$) to a
$15.1$\,ms CLIP ViT-B/16 pass. Training details are given in Section~\ref{sec:exp}.
